
\newgeometry{margin=0.4in}
\begin{titlepage}
	
	\begin{minipage}{0.2\textwidth}
		\raggedleft\includegraphics[scale=0.2]{Title Page/university.png}
	\end{minipage}
	\hfill
	\begin{minipage}{0.5\textwidth}
		\centering
		\ar{وزارة التعليم العالي والبحث العلمي}\\
		\ar{جامعة البصرة}\\
		\ar{كلية التربية للعلوم الصرفة}\\
		\ar{قسم الرياضيات}
	\end{minipage}
	\hfill
	\begin{minipage}{0.2\textwidth}
		\raggedright\includegraphics[scale=0.16]{Title Page/college.png}   
	\end{minipage}
	
	\vspace{1cm}
	
	\begin{center}
		\fontsize{25pt}{30pt}\selectfont
		\vspace{11pt}
				\textbf{\ar{قيود المعاملات لعوائل الدوال متعددة التكافؤ}}\\
	\end{center}
	
	\vfill
	
	\begin{center}
		\large
		\textbf{\ar{مقدم إلى قسم الرياضيات، كلية التربية للعلوم الصرفة، جامعة البصرة\\
				\vspace{6pt}
				كجزء من متطلبات نيل درجة البكالوريوس في علوم الرياضيات}}
	\end{center}
	\vfill
	
	\begin{center}
		\large
		\textbf{\ar{من إعداد الطالبة}}\\
		\vspace{8pt}
		\large
		\textbf{\ar{زهراء عبدالقادر}}
	\end{center}
	
	\vspace{10pt}
	
	\begin{center}
		\large
		\textbf{\ar{بإشراف}}\\
		\vspace{8pt}
		\large
		\textbf{\ar{أ.م.د صارم حازم هادي}}
	\end{center}
	
	\vspace{10pt}
	
	\noindent
	\textbf{\ar{1447 هـ}}
	\hfill
	\textbf{2026 م}
\end{titlepage}
\restoregeometry